\section{\texorpdfstring{射影 $\mathcal{I}$ 算子}{射影 I 算子}}\label{sec:projective-ideal-operator}

射影 $\mathcal{V}$ 算子使我们从多项式集合转到射影空间的子集. 现在转向射影 $\mathcal{I}$ 算子, 它沿相反的方向作用. $\mathcal{I}_{\mathbb{P}}$ 的定义正如所料.

\begin{Definition}{射影 $\mathcal{I}$ 算子}\leavevmode\label{def:projective-ideal-operator}
设 $X\subseteq\mathbb{P}^{n}$ 为子集. $X$ 的\emph{消没理想}是
\[
  \mathcal{I}_{\mathbb{P}}(X)=
    \{f\in K[x_{0},\ldots,x_{n}]\mid f(a)=0\text{ 对所有 }a\in X\}.
\]
若 $K[x_{0},\ldots,x_{n}]$ 的子集对某个 $X\subseteq\mathbb{P}^{n}$ 形如 $\mathcal{I}_{\mathbb{P}}(X)$, 则称它为\emph{射影消没理想}.
\end{Definition}

\begin{BookFlowNote}[l]{%
回顾 Definition~\ref{def:projective-vanishing}, 当我们说 $f(a)=0$ 时, 是指 $f$ 在点 $a\in\mathbb{P}^{n}$ 的每个代表元上取值都为零. 例如, 若 $[1:0]$ 是 $X$ 的一点, 那么多项式 $f=x_{0}-1$ 不属于 $\mathcal{I}(X)$.
}
\noindent 与 $\mathcal{V}_{\mathbb{P}}$ 算子一样, 当由语境可清楚看出是在仿射空间还是射影空间中讨论时, 常写作 $\mathcal{I}(X)$.

尽管 $f(1,0)=0$, 但 $f$ 在代表元 $[2:0]=[1:0]$ 上取值不为零. 下一个例子进一步说明这一点.
\end{BookFlowNote}

\begin{Example}{$\mathbb{P}^{1}$ 中一点的消没理想}\leavevmode\label{ex:vanishing-ideal-of-point-in-projective-line}
令 $X=\{[1:0]\}\subseteq\mathbb{P}^{1}$. 那么 $f=x_{1}\in K[x_{0},x_{1}]$ 是 $\mathcal{I}(X)$ 的元素, 因为点 $[1:0]\in X$ 的任意代表元对某个 $a$ 都形如 $[a:0]$, 而对任意 $a$ 都有 $f(a,0)=0$. 更一般地, 有 $\langle x_{1}\rangle\subseteq\mathcal{I}(X)$. 实际上我们断言等号成立: $\mathcal{I}(X)=\langle x_{1}\rangle$.

为证明反向包含, 令 $f\in\mathcal{I}(X)$. 那么对所有非零 $a\in K$, 有 $f(a,0)=0$. 因此, $f(x_{0},0)$ 是有无穷多个零点的一元多项式, 必为零多项式. 把 $f$ 写成 $(K[x_{0}])[x_{1}]$ 的元素:
\[
  f=\sum_{d\geq0}f_{d}(x_{0})x_{1}^{d}.
\]
由 $f_{0}(x_{0})=f(x_{0},0)=0$, 可得
\[
  f=x_{1}\sum_{d\geq1}f_{d}(x_{0})x_{1}^{d-1}\in\langle x_{1}\rangle.
\]
\end{Example}

\begin{Example}{$\mathbb{P}^{2}$ 中一条直线的消没理想}\leavevmode\label{ex:vanishing-ideal-of-projective-line-in-plane}
若 $X$ 是 Example~\ref{ex:projective-line-in-projective-plane} 中的直线 $\mathcal{V}(x_{0}+x_{1}-x_{2})\subseteq\mathbb{P}^{2}$, 则
\[
  \mathcal{I}(X)=\langle x_{0}+x_{1}-x_{2}\rangle.
\]
$\langle x_{0}+x_{1}-x_{2}\rangle$ 中每个元素在 $X$ 的每一点处取值都为零, 这基本上是直接的. 而反向包含是 Exercise~\ref{ex:vanishing-ideal-of-projective-line-by-substitution} 的内容.
\end{Example}

与仿射情形一样, 射影消没理想确实是理想, 并且容易看出它们是根理想. 不过, 在射影情形还能进一步得到: Lemma~\ref{lem:projective-vanishing-of-components} 蕴含, 对每个 $f\in\mathcal{I}(X)$, $f$ 的每个齐次分量也必属于 $\mathcal{I}(X)$. $\mathcal{I}(X)$ 的这一性质正是齐次理想的定义性质.

\begin{Definition}{齐次理想}\leavevmode\label{def:homogeneous-ideal}
若对理想 $I\subseteq K[x_{0},\ldots,x_{n}]$ 中每个 $f\in I$, $f$ 的每个齐次分量也属于 $I$, 则称 $I$ 是\emph{齐次理想}.
\end{Definition}

上述齐次理想的定义直接来自对消没理想的讨论, 而下面的结果给出齐次理想的另一种重要刻画. 这种刻画或许更直接, 在实践中也很有用.

\begin{Proposition}{齐次理想的刻画}\leavevmode\label{prop:homogeneous-ideal-has-homogeneous-generators}
理想 $I\subseteq K[x_{0},\ldots,x_{n}]$ 是齐次的当且仅当它有一组齐次生成元.
\end{Proposition}
\begin{Proof}
先设 $I$ 是齐次理想, 并令 $S\subseteq I$ 为 $I$ 中所有齐次多项式组成的子集. 只须证明 $I=\langle S\rangle$. 由于 $S\subseteq I$ 且 $I$ 是理想, 有 $\langle S\rangle\subseteq I$. 为证明反向包含, 设 $f\in I$. 那么可将 $f$ 写成非零齐次分量 $f_{k}$ 之和. $I$ 是齐次理想, 所以每个 $f_{k}\in I$. 又由于 $f_{k}$ 齐次, 可知 $f_{k}\in S$. 因此 $f$ 是 $S$ 中元素之和, 即 $f\in\langle S\rangle$. 从而 $I=\langle S\rangle$, 如所断言.

反过来, 设 $I=\langle S\rangle$, 其中 $S$ 是一组齐次多项式. 为证明 $I$ 为齐次理想, 令 $f\in I$. 由 $I=\langle S\rangle$, 可知
\[
  f=\sum_{i=1}^{m}g_{i}h_{i}
\]
其中 $g_{i}\in K[x_{0},\ldots,x_{n}]$ 且 $h_{i}\in S$. 特别地, $h_{i}$ 是某个次数 $d_{i}$ 的齐次多项式. 对每个 $k$, 有
\[
  f_{k}=\sum_{i=1}^{m}(g_{i}h_{i})_{k},
\]
其中右侧的 $(g_{i}h_{i})_{k}$ 表示多项式 $g_{i}h_{i}$ 的 $k$ 次齐次分量. 由 Exercise~\ref{ex:homogeneous-component-of-product-with-homogeneous-factor}, 可以将各加项写成
\[
  (g_{i}h_{i})_{k}=
  \begin{cases}
    (g_{i})_{k-d_{i}}\cdot h_{i} & \text{若 }d_{i}\leq k,\\
    0 & \text{若 }d_{i}>k,
  \end{cases}
\]
其中 $(g_{i})_{k-d_{i}}$ 是 $g_{i}$ 的 $(k-d_{i})$ 次齐次分量. 这蕴含 $(g_{i}h_{i})_{k}\in I$, 因为它是 $h_{i}\in S$ 的倍式. 所以 $f_{k}$ 是 $I$ 中元素之和, 从而 $f_{k}\in I$. 这就证明 $I$ 是齐次理想.
\end{Proof}

理解了齐次理想后, 现在回到本节的主要话题: 消没理想. 下一个结果概括了射影消没理想最重要的代数性质.

\begin{Proposition}{$\mathcal{I}_{\mathbb{P}}(X)$ 是齐次根理想}\leavevmode\label{prop:projective-vanishing-ideal-is-homogeneous-radical}
若 $X\subseteq\mathbb{P}^{n}$ 是任意子集, 则 $\mathcal{I}_{\mathbb{P}}(X)\subseteq K[x_{0},\ldots,x_{n}]$ 是齐次根理想.
\end{Proposition}
\begin{Proof}
$\mathcal{I}_{\mathbb{P}}(X)$ 是根理想的事实, 可以用与仿射情形 (Proposition~\ref{prop:vanishing-ideals-are-radical}) 完全相同的论证得出, 建议读者自行验证. $\mathcal{I}(X)$ 是齐次的, 则按定义由 Lemma~\ref{lem:projective-vanishing-of-components} 得出.
\end{Proof}

和仿射情形一样, 定义消没理想的主要理由是, 为每个射影簇选出一个特别的定义理想. 下一个结果的第三项验证了消没理想的这一作用. 它只是联系射影 $\mathcal{V}$ 和 $\mathcal{I}$ 算子的若干重要性质之一.

\begin{Proposition}{$\mathcal{V}_{\mathbb{P}}$ 和 $\mathcal{I}_{\mathbb{P}}$ 的基本性质}\leavevmode\label{prop:basic-properties-of-projective-vanishing-operators}
设 $S,T\subseteq K[x_{0},\ldots,x_{n}]$ 和 $X,Y\subseteq\mathbb{P}^{n}$ 为子集.
\begin{enumerate}
  \item 若 $S\subseteq T$, 则 $\mathcal{V}_{\mathbb{P}}(S)\supseteq\mathcal{V}_{\mathbb{P}}(T)$.
  \item 若 $X\subseteq Y$, 则 $\mathcal{I}_{\mathbb{P}}(X)\supseteq\mathcal{I}_{\mathbb{P}}(Y)$.
  \item $\mathcal{V}_{\mathbb{P}}(\mathcal{I}_{\mathbb{P}}(X))\supseteq X$, 等号成立当且仅当 $X$ 为射影簇.
  \item $\mathcal{I}_{\mathbb{P}}(\mathcal{V}_{\mathbb{P}}(S))\supseteq S$, 等号成立当且仅当 $S$ 为射影消没理想.
\end{enumerate}
\end{Proposition}
\begin{Proof}
这些陈述的证明与仿射情形中相应陈述的证明类似 (Propositions~\ref{prop:operators-reverse-inclusions} 和~\ref{prop:compose-vanishing-operators}), 建议读者自行验证.
\end{Proof}

继续与仿射情形作比较. 回忆, 在仿射情形, 我们利用 $\mathcal{V}$ 和 $\mathcal{I}$ 算子间的关系证明了不可约分解的存在性和唯一性. 现在陈述射影情形中相应的结论, 从不可约射影簇的自然定义开始. 这个定义从仿射情形原样沿用.

\begin{Definition}{不可约射影簇}\leavevmode\label{def:irreducible-projective-variety}
若射影簇 $X\subseteq\mathbb{P}^{n}$ 对某些射影簇 $X_{1},X_{2}\subsetneq X$ 满足 $X=X_{1}\cup X_{2}$, 则称 $X$ \emph{可约}. 若 $X$ 既不是空集, 也不是可约簇, 则称 $X$ \emph{不可约}.
\end{Definition}

正如可以预料的那样, 在射影情形, 不可约分解总是存在且唯一, 证明与仿射情形类似.

\begin{BookPropositionDefinition}{不可约分解}\leavevmode\label{prop:projective-irreducible-decomposition}
设 $X\subseteq\mathbb{P}^{n}$ 为非空射影簇. 则存在不可约射影簇 $X_{1},\ldots,X_{r}\subseteq X$, 使得对任意 $i\neq j$ 有 $X_{i}\not\subseteq X_{j}$, 且
\begin{equation}\label{eq:decomposition-into-projective-irreducible-components}
  X=\bigcup_{i=1}^{r}X_{i}.
\end{equation}
此外, 在不计排列顺序的意义下, 这些不可约射影簇 $X_{1},\ldots,X_{r}$ 是唯一的. 称它们为 $X$ 的\emph{不可约分支}, 称 \eqref{eq:decomposition-into-projective-irreducible-components} 为 $X$ 的\emph{不可约分解}.
\end{BookPropositionDefinition}
\begin{Proof}
证明用到射影 $\mathcal{V}$ 和 $\mathcal{I}$ 算子间的关系, 以及 $K[x_{0},\ldots,x_{n}]$ 的 Noether 性质, 与仿射情形的陈述 (Proposition/Definition~\ref{prop:irreducible-decomposition}) 的证明类似. 把验证留给读者作为练习.
\end{Proof}

现在, 射影簇和仿射簇在结构上的相似性已经开始显现. 射影情形中的不少结果, 其证明确实与仿射情形相同或至少类似. 然而, 我们对射影簇还缺乏几何直观. 毕竟, 即使在实数域上, $n\geq2$ 时也很难想象 $\mathbb{P}^{n}$ 的样子, 又如何画出射影簇的图形? 回答这个问题的关键是两种在射影情形和仿射情形之间转换的方法: 仿射限制与射影闭包. 接下来两节专门讨论这两种方法.

\BookExerciseHeading
\begin{BookExercises}
  \item\label{ex:vanishing-ideal-of-projective-line-by-substitution} 令 $X=\mathcal{V}(x_{0}+x_{1}-x_{2})\subseteq\mathbb{P}^{2}$. 我们证明 $\mathcal{I}(X)=\langle x_{0}+x_{1}-x_{2}\rangle$. 如 Example~\ref{ex:vanishing-ideal-of-projective-line-in-plane} 所述, 只须证明 $\mathcal{I}(X)\subseteq\langle x_{0}+x_{1}-x_{2}\rangle$.
  \begin{enumerate}
    \item 令 $f\in\mathcal{I}(X)$. 证明 $f(x_{0},x_{1},x_{0}+x_{1})=0\in K[x_{0},x_{1}]$.
    \item 论证, 由 $f(x_{0},x_{1},x_{0}+x_{1})=0$ 可得 $f\in\langle x_{0}+x_{1}-x_{2}\rangle$.
  \end{enumerate}
  \item\label{ex:homogeneous-component-of-product-with-homogeneous-factor} 令 $g\in K[x_{0},\ldots,x_{n}]$ 为任意多项式, 且令 $h\in[x_{0},\ldots,x_{n}]$ 为 $d$ 次齐次多项式. 证明对每个 $k\geq0$, 有
  \[
    (gh)_{k}=
    \begin{cases}
      g_{k-d}\cdot h & \text{若 }d\leq k,\\
      0 & \text{若 }d>k,
    \end{cases}
  \]
  其中 $(gh)_{k}$ 是 $gh$ 的 $k$ 次齐次分量, 而 $g_{k-d}$ 是 $g$ 的 $(k-d)$ 次齐次分量.
  \item\label{ex:identify-homogeneous-ideals-in-two-variables} $K[x,y]$ 中下列哪些理想是齐次的?
  \begin{enumerate}
    \item $\langle x+1,y^{2}\rangle$
    \item $\langle x+y,x^{2}\rangle$
    \item $\langle x^{2}+y,x^{2}-y\rangle$
    \item $\langle x^{2}+y^{2},xy^{2}+y^{3}+x^{2},y^{2}-x^{2}\rangle$
  \end{enumerate}
  \item\label{ex:principal-ideal-homogeneous-iff-generator-homogeneous} 证明主理想 $\langle f\rangle\subseteq K[x_{0},\ldots,x_{n}]$ 是齐次的当且仅当多项式 $f$ 是齐次的.
  \item\label{ex:projective-irreducibility-equivalent-to-prime-ideal} 设 $X\subseteq\mathbb{P}^{n}$ 为射影簇. 证明 $X$ 不可约当且仅当 $\mathcal{I}_{\mathbb{P}}(X)$ 为素理想.
  \item\label{ex:projective-single-point-equivalent-to-maximal-ideal} 证明子集 $X\subseteq\mathbb{P}^{n}$ 只含一个点当且仅当 $\mathcal{I}_{\mathbb{P}}(X)$ 为极大理想.
  \item\label{ex:adapt-affine-proof-of-radical-ideal-to-projective-space} 用射影簇的语言和记号改写 Proposition~\ref{prop:vanishing-ideals-are-radical} 的证明, 确保每一步都能在射影情形下完成, 从而证明射影消没理想是根理想.
  \item\label{ex:adapt-affine-proof-of-operator-properties-to-projective-space} 用射影簇的语言和记号改写 Propositions~\ref{prop:operators-reverse-inclusions} 和~\ref{prop:compose-vanishing-operators} 的证明, 确保每一步都能在射影情形下完成, 从而证明 Proposition~\ref{prop:basic-properties-of-projective-vanishing-operators}.
  \item\label{ex:adapt-affine-proof-of-decomposition-to-projective-space} 用射影簇的语言和记号改写 Proposition/Definition~\ref{prop:irreducible-decomposition} 的证明, 确保每一步都能在射影情形下完成, 从而证明 Proposition/Definition~\ref{prop:projective-irreducible-decomposition}.
\end{BookExercises}
